% ============================================================================
\chapter{Installation}
\label{ch:installation}
% ============================================================================

This chapter covers building the application from source and running it on a
development machine. Deployment on the Raspberry~Pi operator terminal is the
subject of Chapter~\ref{ch:rpi}.

\section{Prerequisites}

\begin{itemize}[itemsep=2pt]
  \item \textbf{Rust toolchain} --- install via \url{https://rustup.rs}. The project
        uses Rust edition 2024; any recent stable toolchain works.
  \item \textbf{Build dependencies} --- a C compiler and the graphics/input
        development libraries required by Slint. On Debian-based systems:
\begin{lstlisting}[language=bash]
$ sudo apt install build-essential cmake libfontconfig1-dev \
      libfreetype6-dev libxkbcommon-dev libinput-dev \
      libudev-dev libgbm-dev
\end{lstlisting}
  \item \textbf{Serial port access} --- the user running the application must be able
        to open the robot's serial device:
\begin{lstlisting}[language=bash]
$ sudo usermod -aG dialout $USER    # log out and back in afterwards
\end{lstlisting}
\end{itemize}

\section{Building from source}

\begin{lstlisting}[language=bash]
$ git clone <repository-url> deltax2
$ cd deltax2
$ cargo build --release
\end{lstlisting}

The binary is produced at \code{target/release/deltax2}.

\begin{note}
The application reads \code{config.toml} from its \emph{current working directory}
at start-up, not from the directory containing the binary. Always start it from a
directory that contains a valid \code{config.toml} (this matters for the systemd
unit in Section~\ref{sec:systemd}, which sets \code{WorkingDirectory} explicitly).
\end{note}

\section{Desktop test run}

For development on a normal Linux desktop, connect the robot (or leave it
disconnected --- the UI starts regardless and shows \emph{Disconnected}), then:

\begin{lstlisting}[language=bash]
$ cargo run
\end{lstlisting}

Configure the serial port and tray layouts as described in
Chapter~\ref{ch:configuration} before first use with real hardware.
